\documentclass[11pt,reqno]{amsart}

%==============================================================
% Iterated Brownian motion does not preserve the
% Laguerre--Polya class
%==============================================================

\usepackage[utf8]{inputenc}
\usepackage[T1]{fontenc}
\usepackage{lmodern}
\usepackage[margin=1.15in]{geometry}
\usepackage{amsmath,amssymb,amsthm,mathtools}
\usepackage{microtype}
\usepackage{hyperref}
\emergencystretch=3em

\theoremstyle{plain}
\newtheorem{theorem}{Theorem}[section]
\newtheorem{proposition}[theorem]{Proposition}
\newtheorem{lemma}[theorem]{Lemma}
\newtheorem{corollary}[theorem]{Corollary}
\newtheorem{conjecture}[theorem]{Conjecture}
\theoremstyle{definition}
\newtheorem{definition}[theorem]{Definition}
\theoremstyle{remark}
\newtheorem{remark}[theorem]{Remark}
\newenvironment{keyresult}{\par\medskip\noindent}{\par\medskip}

\newcommand{\R}{\mathbb R}
\newcommand{\C}{\mathbb C}
\newcommand{\N}{\mathbb N}
\newcommand{\E}{\mathbb E}
\newcommand{\LP}{\mathcal{LP}}
\DeclareMathOperator{\sgn}{sgn}

\title[Nonmonotone Laguerre--P\'olya dynamics for the quartic flow]{Nonmonotone
Laguerre--P\'olya Dynamics for the Quartic de Bruijn Flow}
\author{Ramiro Fontes}
\address{Quijotic Research}
\email{ramirofontes@gmail.com}
\urladdr{www.quijotic.com}
\date{\today}

\begin{document}

\begin{abstract}
The growing multiplier $e^{tu^2}$, the backward heat operator
$e^{-t\partial_x^2}$, preserves the Laguerre--P\'olya class (P\'olya), the
analytic engine of de Bruijn--Newman theory. We study its order-four analogue,
the biharmonic pseudo-process semigroup $e^{t\partial_x^4}$, from the signed,
two-sided iterated-Brownian constructions of Funaki and Allouba--Zheng, multiplier
$e^{tu^4}$, and show it does not preserve the
class. The counterexample is a quartic polynomial,
\[
   e^{t\partial_x^4}\,(x^2-1)(x^2-4)\;=\;x^4-5x^2+(4+24t),
\]
real-rooted for $t\le 3/32$ and not after; the failure is sharp in the order,
the multiplier $e^{tu^{2m}}$ preserving $\LP$ if and only if $m=1$. At the de
Bruijn--Newman boundary the same lifting reaches the Riemann $\Xi$-function: at
$t=8$, $H_{8,4}=e^{8\partial_x^4}H_0$ violates the Laguerre inequality at an
explicit point, so $H_{8,4}\notin\LP$, a theorem certified in validated
interval arithmetic. Since $H_0=\tfrac18\Xi(\cdot/2)$ lies in $\LP$ if and only if
the Riemann hypothesis holds, the theorem itself is the membership failure
$H_{8,4}\notin\LP$; should RH hold, so that $\Xi$ began in $\LP$, it says the
quartic flow has carried $\Xi$ out of the class. A two-saddle steepest-descent estimate, carried into the
complex plane and closed by Rouch\'e, then real-roots $H_{t,4}$ on a window
$|z|\le X_t\to\infty$ that widens past the saddle band, unlike the order-two flow,
which once inside $\LP$ stays. Left open is the far field $|x|>X_t$, and with it
global return $\Lambda(4)<\infty$; the window theorem is asymptotic in $t$.
\end{abstract}

\maketitle

%==============================================================
\section{Introduction}\label{sec:intro}
%==============================================================

A real entire function belongs to the \emph{Laguerre--P\'olya class} $\LP$ if
it is a locally uniform limit of real polynomials with only real zeros;
equivalently, by the Hadamard factorization,
$f(x)=cx^m e^{-\gamma x^2+\beta x}\prod_k(1-x/x_k)e^{x/x_k}$ with
$\gamma\ge0$, $c,\beta,x_k\in\R$ and $\sum x_k^{-2}<\infty$~\cite{Levin}.

The closure property at the root of de Bruijn--Newman theory is invariance under
the backward heat operator. On the cosine-transform side, multiplication
by the growing factor $e^{tu^2}$, equivalently the operator $e^{-t\partial_x^2}$
on $x$, preserves the class:
\begin{equation}\label{eq:heatLP}
   f\in\LP\ \Longrightarrow\ e^{-t\partial_x^2}f\in\LP\qquad(t\ge0).
\end{equation}
This is P\'olya's theorem on universal factors~\cite{Polya13,Polya27,deBruijn},
the multiplier $e^{tu^2}=e^{\gamma w}$ ($w=u^2$, $\gamma=t\ge0$) being of the
admissible form. The sign is essential. The forward Gaussian convolution
$G_\tau f(x)=\E\,f(x+\sqrt{2\tau}\,N)=e^{+\tau\partial_x^2}f$, multiplier
$e^{-\tau u^2}$, does not preserve $\LP$: already
$G_\tau(x^2-1)=x^2-1+2\tau$ has non-real zeros for $\tau>\tfrac12$. (Forward
convolution is variation-diminishing in Schoenberg's sense, but a
variation-diminishing map can still send real zeros into the complex plane, as
this example does; reduction of sign changes is not preservation of
real-rootedness.) It is the backward direction $e^{tu^2}$ that lets de Bruijn and
Newman run the flow $H_t(x)=\int_0^\infty e^{tu^2}\Phi(u)\cos(xu)\,du$ on the
Riemann $\Xi$-function ($\partial_t H_t=-\partial_x^2H_t$) and speak of a
sharp real-rooting threshold $\Lambda$, the de Bruijn--Newman
constant~\cite{deBruijn,Newman,RodgersTao}.

The de Bruijn multiplier is one of a one-parameter family: for $\alpha>0$,
\begin{equation}\label{eq:family}
   H_{t,\alpha}(x)=\int_0^\infty e^{tu^\alpha}\Phi(u)\cos(xu)\,du,
\end{equation}
with $\alpha=2$ the flow above. Among the symbols $e^{\lambda u^\alpha}$ only
$\alpha=2$ is a P\'olya universal
factor, so only there does the flow preserve $\LP$; and for $\alpha\notin2\N$ the
corner of $u^\alpha$ at the origin destroys the same entire-function setting, so
the even exponents are the natural real-rooting deformations. The first even exponent past the classical
$\alpha=2$ is $\alpha=4$, the quartic, biharmonic, iterated-Brownian flow: the
first nonclassical member of the family, and the subject of this paper.

On $x$ this is the signed biharmonic pseudo-process semigroup
\begin{equation}\label{eq:P4}
   P_t=e^{t\partial_x^4}\qquad(\text{multiplier }e^{tu^4}),
\end{equation}
the order-four object of the signed, two-sided Brownian-time and
iterated-Brownian constructions of Funaki and Burdzy~\cite{Funaki,Burdzy} (a
Brownian motion run along an independent Brownian clock). The kernel of $P_t$ is
signed, so $P_t$ is a pseudo-process semigroup, not the Markov semigroup of a
process (\S\ref{sec:ibm}). The Gaussian-clock identity
$e^{tu^4}=\E_{s\sim N(0,2t)}\,e^{su^2}$ realizes $P_t$ as the de Bruijn flow
$e^{su^2}$ run for a Gaussian-random, two-sided time $s$, the sense in which it
iterates the Brownian and the de Bruijn evolutions. Allouba and Zheng~\cite{AlloubaZheng}
placed these Brownian-time processes on a rigorous footing and identified the
half-derivative generator and the biharmonic operator behind them; Nane and
Baeumer--Meerschaert--Nane developed the associated higher-order and fractional
Cauchy problems~\cite{Nane,BMN}; the signed transition kernels go back to Krylov
and Hochberg~\cite{Krylov,Hochberg}.

One may ask whether the preservation \eqref{eq:heatLP} survives the
passage to order four: does $P_t$ preserve $\LP$? The de Bruijn flow $e^{tu^2}$
($m=1$) does; the order-four flow $e^{tu^4}$ ($m=2$) is the next case, and the
canonical probabilistic evolution past Brownian motion. De
Bruijn~\cite{deBruijn} proved that the heat flow brings the Riemann
$\Xi$-function, of order one and maximal type, into the Laguerre--P\'olya
class in finite time, $H_t\in\LP$ for $t\ge\tfrac12$. The iterate behaves
oppositely: $H_{t,4}\notin\LP$ at some finite $t$ (certified at $t=8$,
\S\ref{sec:zeta}), and if RH holds, so that $\Xi$ began inside $\LP$, this is the
iterate carrying $\Xi$ out of the class. A two-saddle steepest-descent estimate,
taken into the complex plane and closed by Rouch\'e, then proves that for large
$t$ the flow re-roots on a window that widens without bound: heat enters $\LP$
and stays, the order-four flow leaves and re-roots on the bulk, the far field
the one piece still open.

\subsection*{Results}
Four results. $P_t$ fails to preserve $\LP$ already on a polynomial; the failure
is sharp in the order; the boundary flow $H_{t,4}$ leaves $\LP$ at finite time;
and for large $t$ it re-roots on a window that widens past the saddle band. The
first:

\begin{keyresult}
\textbf{Theorem~\ref{thm:witness}.} $P_t\,(x^2-1)(x^2-4)=x^4-5x^2+(4+24t)$ has
four non-real zeros, two conjugate pairs, for every $t>3/32$. Hence $P_t$ does
not preserve $\LP$.
\end{keyresult}

On a polynomial $P_t=e^{t\partial_x^4}$ is the finite sum $f+tf''''+\cdots$, so
the constant floor $24t$ fills the dip of $f$ and lifts the pair off the axis at
$t=3/32$ (\S\ref{sec:nonpres}).

For polynomials, non-preservation also follows from the Borcea--Br\"and\'en
classification of linear preservers of real-rootedness~\cite{BorceaBranden} with
the P\'olya--Schur theory~\cite{PolyaSchur}: a diagonal symbol preserves
real-rootedness only if it is itself of Laguerre--P\'olya type, and
$z\mapsto e^{tz^4}$ is not. We give instead the explicit counterexample, a probabilistic
reading of the evolution (Proposition~\ref{prop:mixture}, Lemma~\ref{lem:mixture}),
a pointwise Laguerre certificate (Lemma~\ref{lem:laguerre}), and the application
to $\Xi$ in \S\ref{sec:zeta}.

By the identity $e^{tu^4}=\E_{s\sim N(0,2t)}e^{su^2}$, the evolution is the de
Bruijn flow run over a two-sided Gaussian time (\S\ref{sec:ibm}),
\begin{equation}\label{eq:mix0}
   P_tf=\E_{s\sim N(0,2t)}\bigl[\,e^{-s\partial_x^2}f\,\bigr];
\end{equation}
a convex combination of two $\LP$ functions can leave the class, already
$\tfrac12[(x-1)(x-2)+(x-3)(x-4)]=x^2-5x+7$ (Lemma~\ref{lem:mixture}). The analogy
is imperfect, since for $s<0$ the summand $e^{|s|\partial_x^2}f$ is the forward
Gaussian and need not lie in $\LP$; the proof is instead the expansion
$P_tf=f+tf''''+\cdots$ on the counterexample, with $f'^2-ff''<0$ at a point
(Lemma~\ref{lem:laguerre}). By Theorem~\ref{thm:order}, $e^{tu^{2m}}$ preserves
$\LP$ if and only if $m=1$.

Funaki's representation makes $P_t$ act on the Riemann $\Xi$-function as the
quartic de Bruijn--Newman flow
$H_{t,4}(x)=\int_0^\infty e^{tu^4}\Phi(u)\cos(xu)\,du=P_tH_0(x)$, where
$H_0(x)=\tfrac18\Xi(x/2)$ is the classical Fourier representation of the
completed zeta function~\cite{Titchmarsh,Polya27}. Section~\ref{sec:zeta} shows it first leaves the class:

\begin{keyresult}
\textbf{Proposition~\ref{prop:leave}.} $H_{8,4}\notin\LP$: validated interval
arithmetic (Arb) gives $L_1[H_{8,4}](63.2083)\in[-2.3966,-2.3965]\times10^{-20}<0$.
If RH holds, so that $\Xi$ began in $\LP$, this says the quartic flow has carried
$\Xi$ out of the class.
\end{keyresult}

A two-saddle analysis then shows it real-rooted on the bulk of the spectrum:

\begin{keyresult}
\textbf{Theorem~\ref{thm:window}.} For all large $t$, every zero of $H_{t,4}$ in
the disk $|z|\le X_t$, $X_t\to\infty$, is real and simple. The proof is a uniform
two-saddle approximation of $H_{t,4}$ closed by a Rouch\'e argument.
\end{keyresult}

Whether the real-rooting reaches every height, not only the widening window, is
the question left open:

\begin{keyresult}
\textbf{Conjecture~\ref{conj:finite}.} $\Lambda(4)<\infty$: for large $t$,
$H_{t,4}$ is real-rooted at every height; equivalently, the exit is transient.
\end{keyresult}

The departure is arithmetic. The pair $\{2\gamma_4,2\gamma_5\}$ brackets a
positive hump of $H_0$, and the order-four floor there, the value of $H_0''''$
at the midpoint, is negative, so the iterated-Brownian average pushes the hump
beneath the axis and the two real zeros collide into a conjugate pair, a
Laguerre violation certified at a single point. Unusually close zero pairs are the input the order-four floor converts into a
non-real pair (in the extreme, the anomalously close pairs of the Lehmer
phenomenon~\cite{CSV} at large heights); the only arithmetic datum used here is
the location of
two low-lying critical-line zeros. Re-entry, by contrast, is not arithmetic at
all: for large $t$ the
weight $e^{tu^4}$ concentrates the transform at an interior saddle
$u_*\approx\tfrac14\log t$, overwriting the arithmetic of $H_0$ with a
Gaussian-damped cosine whose leading term is real-rooted. Theorem~\ref{thm:window}
turns this into real-rooting on the widening window $|z|\le X_t$; carrying the
same two-saddle dominance into the far field $|x|>X_t$ is the one estimate we
leave open. The same is conjectured at every even order $2m\ge4$
(\S\ref{sec:zeta}), with order six
supplying supporting computations: the multiplier $e^{tu^{2m}}$ is non-universal,
the flow $H_{t,2m}$ fails the Laguerre certificate at finite $t$ (certified for
$m=2$, computed for $m=3$), and the kernel forces the saddle to detach, so it is
conjecturally regained, $\Lambda(2m)<\infty$. This would make order two
the unique even order under which $\LP$ is forward-invariant, a sharpening of
Theorem~\ref{thm:order} from preservation to permanence.

The $\Xi$-function enters only as the boundary example where \eqref{eq:heatLP}
meets the de Bruijn--Newman theory; the inputs are classical.

%==============================================================
\section{The biharmonic/iterated-Brownian semigroup}\label{sec:ibm}
%==============================================================

We construct $P_t$ as the semigroup of Funaki's iterated Brownian motion, with
multiplier $e^{tu^4}$, and rewrite it as the de Bruijn flow averaged over a
Gaussian time.

\subsection{Construction}
Let $W$ and $V$ be independent two-sided standard Brownian motions. Funaki's
iterated Brownian motion is
\begin{equation}\label{eq:Z}
   Z_s=\begin{cases} W(V_s), & V_s\ge0,\\[2pt]
   i\,\widetilde W(-V_s), & V_s<0,\end{cases}
\end{equation}
with $\widetilde W$ a third independent Brownian motion and the imaginary
branch read through the entire extension of the test
function~\cite{Funaki}. For $f$ entire of order at most one and maximal type,
the expectation $\E\,f(x+Z_{8t})$ converges, and conditioning on the clock $V$
gives, for any such $f$,
\begin{equation}\label{eq:cond}
   \E\bigl[f(x+Z_{8t})\mid V\bigr]=G_{V_{8t}/2}\,f(x),
   \qquad \tfrac12 V_{8t}\sim N(0,2t),
\end{equation}
where $G_\tau f(x)=\E\,f(x+\sqrt{2\tau}\,N)$ for $\tau\ge0$ is the heat (Gauss)
average and $G_\tau$ for $\tau<0$ is its backward continuation, the de
Bruijn--Newman multiplier $e^{-\tau u^2}$ on the Fourier side. The verification
of \eqref{eq:cond} is the one-line Gaussian computation: for $\sigma\ge0$,
$\E\,f(x+W_\sigma)=G_{\sigma/2}f(x)$, and for the imaginary branch
$\E\,f(x+i\widetilde W_\sigma)=G_{-\sigma/2}f(x)$ since the odd part of $f$
averages out over the symmetric Gaussian and $\E\cosh(\widetilde W_\sigma u)
=e^{\sigma u^2/2}$.

\begin{proposition}[The biharmonic semigroup]\label{prop:semigroup}
For $f$ entire of order $\le1$, maximal type, with $\int_0^\infty
e^{tu^4}|\hat f(u)|\,du<\infty$,
\[
   P_tf(x):=\E\,f(x+Z_{8t})
\]
is well defined, solves $\partial_tP_tf=\partial_x^4P_tf$ with $P_0f=f$, and
acts on the Fourier side by the multiplier $e^{tu^4}$. In particular
$(P_t)_{t\ge0}$ is a semigroup, $P_t=e^{t\partial_x^4}$, and on polynomials
$P_tf=\sum_{k\ge0}\frac{t^k}{k!}f^{(4k)}$, a finite sum.
\end{proposition}

\begin{proof}
By \eqref{eq:cond} and the tower property, $P_tf=\E_V[G_{V_{8t}/2}f]$. Writing
$s=-\tfrac12V_{8t}\sim N(0,2t)$, so that $H_sf:=e^{-s\partial_x^2}f=G_{-s}f=
G_{V_{8t}/2}f$ has Fourier multiplier $e^{su^2}$, gives
$P_tf=\E_{s\sim N(0,2t)}[H_sf]$. Averaging the multiplier
over the centered Gaussian clock,
$\E_{s\sim N(0,2t)}e^{su^2}=e^{(2t)u^4/2}=e^{tu^4}$, the Gaussian moment
generating function evaluated at $u^2$; thus $\widehat{P_tf}=e^{tu^4}\hat f$,
which the hypothesis $\int e^{tu^4}|\hat f|<\infty$ makes a well-defined entire
function, and which gives the semigroup property and the equation
$\partial_tP_tf=\partial_x^4P_tf$. For a polynomial $f$ the $L^1$ hypothesis is
inapplicable but unnecessary: $P_tf(x)=\E\,f(x+Z_{8t})=\sum_{n\ge0}
\tfrac{1}{n!}\E[Z_{8t}^n]\,f^{(n)}(x)$ is a finite sum, and the moments of the
iterated Brownian motion are $\E Z_{8t}^{2}=0$ and, conditioning on the clock,
$\E Z_{8t}^{4}=\E\bigl[3V_{8t}^2\bigr]=3\,\E V_{8t}^2=24t$, more generally
$\E Z_{8t}^{4k}=\tfrac{(4k)!}{k!}t^k$ and $\E Z_{8t}^{n}=0$ for $4\nmid n$: these
are the moments of $e^{t\partial_x^4}$, so the two definitions agree term by
term and $P_tf=\sum_{k\ge0}\tfrac{t^k}{k!}f^{(4k)}$.
\end{proof}

\begin{remark}[The two-sided clock and the memory term]\label{rem:allouba}
The two-sided clock with the imaginary branch is what produces the clean
generator $\partial_x^4$. The one-sided Brownian-time process $W(|V_s|)$ of
Allouba and Zheng~\cite{AlloubaZheng} solves instead
$\partial_su=\tfrac{\Delta f}{\sqrt{8\pi s}}+\tfrac18\Delta^2u$, carrying a
memory term in the initial data; its clock $|V_s|$ has the law of the inverse
half-stable subordinator, whence the time-fractional reading~\cite{BMN}. Funaki's
two-sided construction deletes the memory term and leaves the pure
biharmonic evolution~\cite{Funaki}, which is the semigroup
Proposition~\ref{prop:semigroup} describes and the one whose action on $\LP$
we study.
\end{remark}

\subsection{The mixture representation}
The content of \eqref{eq:cond}, restated, is the representation that drives
everything below.

\begin{proposition}[Two-sided heat-clock representation]\label{prop:mixture}
For admissible $f$,
\begin{equation}\label{eq:mix}
   P_tf=\E_{s\sim N(0,2t)}\bigl[\,H_sf\,\bigr],
   \qquad H_sf:=e^{-s\partial_x^2}f\ \ (\text{multiplier }e^{su^2}),
\end{equation}
an average over the centered Gaussian clock of the de Bruijn--Newman
heat flows $H_s$, taken over both signs of the time $s$.
\end{proposition}

The clock is two-sided. For $s\ge0$ the de Bruijn flow
$H_s=e^{-s\partial_x^2}$ preserves $\LP$; for $s<0$ it is the forward Gaussian
$e^{|s|\partial_x^2}$, which does not (\S\ref{sec:intro}). The iterated-Brownian
evolution averages the preserving direction against the destroying one, so the
average need not be real-rooted. We treat \eqref{eq:mix} as the heuristic behind
the non-preservation; the rigorous statement is the counterexample of
\S\ref{sec:nonpres}, where for a polynomial this average is the finite sum
$P_tf=f+tf''''+\cdots$.

%==============================================================
\section{Non-preservation}\label{sec:nonpres}
%==============================================================

We prove $P_t$ does not preserve $\LP$: the class is not convex, an explicit
quartic leaves it under $P_t$, a single Laguerre inequality certifies the
departure, and order two is the only even order that preserves $\LP$.

\subsection{Nonconvexity of the Laguerre--P\'olya class}

\begin{lemma}\label{lem:mixture}
The Laguerre--P\'olya class is not convex. Indeed
\[
   \tfrac12\bigl[(x-1)(x-2)+(x-3)(x-4)\bigr]=x^2-5x+7
\]
is an average of two real-rooted polynomials and has discriminant
$25-28=-3<0$.
\end{lemma}

The two summands are real-rooted with disjoint zero sets ($\{1,2\}$ and
$\{3,4\}$); the average splits the difference and the four real zeros collapse
to a conjugate pair near $x=5/2$. A convex combination of $\LP$ functions can
thus leave $\LP$. This is the heuristic for what follows, with the discrete
average replaced by the Gaussian clock \eqref{eq:mix}; but the analogy is
imperfect: that clock also runs the forward-heat direction $s<0$, whose
summands need not be real-rooted, so we establish non-preservation directly,
by the counterexample of \S\ref{sec:nonpres}, rather than from convexity failure.

\subsection{The counterexample}

\begin{theorem}[Iterated Brownian motion does not preserve $\LP$]
\label{thm:witness}
Let $f(x)=(x^2-1)(x^2-4)\in\LP$. Then
\begin{equation}\label{eq:witness}
   P_tf(x)=x^4-5x^2+(4+24t),
\end{equation}
which lies in $\LP$ for $0\le t\le 3/32$ and has four non-real zeros,
two conjugate pairs, for every $t>3/32$. Consequently the iterated-Brownian
semigroup $P_t=e^{t\partial_x^4}$ does not preserve the Laguerre--P\'olya
class.
\end{theorem}

\begin{proof}
Since $f$ is a polynomial of degree four, $f''''\equiv24$ and $f^{(8)}\equiv0$,
so by Proposition~\ref{prop:semigroup} $P_tf=f+tf''''=x^4-5x^2+4+24t$. Writing
$y=x^2$, the zeros solve $y^2-5y+(4+24t)=0$, $y=\tfrac12\bigl(5\pm\sqrt{9-96t}
\bigr)$. For $t\le3/32$ the discriminant $9-96t\ge0$ and both roots $y$ are
real and positive (their sum $5$ and product $4+24t$ are positive), giving four
real $x$; the four zeros are real and $P_tf\in\LP$. At $t=3/32$ the two $y$-roots
coincide at $5/2$, so $x=\pm\sqrt{5/2}$ are double real zeros. For $t>3/32$ the
discriminant is negative, $y$ is a conjugate pair, and the four zeros
$x=\pm\sqrt{y}$ form two conjugate pairs off the real axis. Thus $P_tf\notin\LP$
for $t>3/32$.
\end{proof}

The threshold is the floor-versus-dip balance in transparent form. On
$[1,2]$ the function $f$ dips to its minimum
$f(\sqrt{5/2})=(\tfrac52-1)(\tfrac52-4)=-\tfrac94$; the evolution raises every
value by the constant $24t$; and the dip is erased, the pair $\{1,2\}$ (and
its mirror $\{-1,-2\}$) lifted, when $24t=\tfrac94$, i.e.
$t=\tfrac{9}{96}=\tfrac3{32}$.

\subsection{The certificate}
The constant floor $f''''=24$ is the one simplification here; the general
obstruction is detected pointwise. For $f\in\LP$ the Laguerre
inequality holds at every real point, and a single violation certifies
non-membership, with no zero-tracking and no contour bookkeeping.

\begin{lemma}[Laguerre certificate]\label{lem:laguerre}
Let $f\in\LP$ be real. Then
\[
   L_1[f](x):=f'(x)^2-f(x)\,f''(x)\ \ge\ 0\qquad(x\in\R).
\]
Hence if $L_1[g](x_0)<0$ at a single real $x_0$, then $g\notin\LP$. In
particular any real entire $g$ with a strict positive local minimum
($g(x_0)>0$, $g'(x_0)=0$, $g''(x_0)>0$) fails to lie in $\LP$.
\end{lemma}

\begin{proof}
By the Hadamard factorization of the class,
$f(x)=cx^m e^{-\gamma x^2+\beta x}\prod_k(1-x/x_k)e^{x/x_k}$ with $\gamma\ge0$
and real zeros $x_k$, $\sum_k x_k^{-2}<\infty$. Off the zeros,
\[
   \Bigl(\frac{f'}{f}\Bigr)'=\frac{f''f-(f')^2}{f^2}
   =-2\gamma-\frac{m}{x^2}-\sum_k\frac1{(x-x_k)^2}\ \le\ 0,
\]
so $L_1[f]=(f')^2-f''f\ge0$ off and, by continuity, on all of $\R$. The
contrapositive is the certificate; at a positive local minimum
$L_1[g]=-g\,g''<0$. The inequality is classical
(Laguerre; see Csordas--Varga~\cite{CsordasVarga}).
\end{proof}

A filled dip is a Laguerre violation: when the order-four floor raises the
enclosed minimum of a real pair above zero, that minimum is a strict positive
local minimum of $P_tf$, and $L_1[P_tf]<0$ there. To prove
$P_tf\notin\LP$ one evaluates $P_tf,(P_tf)',(P_tf)''$ at a single point. This
reproves the counterexample at once: for $t>3/32$ the polynomial $x^4-5x^2+(4+24t)$
has positive local minima at $x=\pm\sqrt{5/2}$, with value $24t-\tfrac94>0$ and
second derivative $20>0$. Being pointwise, the certificate removes any need to follow
a configuration through a range of times: each $t$ is certified on its
own, at its own point.

\begin{remark}[The floor and its sign]\label{rem:floor}
The evolution raises $f$ by the floor $P_tf-f=t\,f''''+O(t^2)$
(Proposition~\ref{prop:semigroup}), positive where $f''''>0$. Once the floor clears the dip depth, a real pair of
$f$ bounding a dip lifts: its enclosed minimum crosses zero from below, and
the certificate fires. This is one-sided: where $f''''<0$ the floor deepens the
dip and the pair stays real.
Non-preservation is therefore the existence of a pair over a positive
floor, which the counterexample exhibits with $f''''\equiv24$; the $\Xi$-corollary of
\S\ref{sec:zeta} asks which pairs sit over a positive floor, and whether the
certificate fires along the flow for arbitrarily large $t$.
\end{remark}

\subsection{Sharpness in the order}
Among the growing-multiplier orders $e^{tu^{2m}}$, the order-two de Bruijn flow
alone preserves the class.

\begin{theorem}[Order dichotomy]\label{thm:order}
For integer $m\ge1$ let $P_{m,t}:=e^{(-1)^m t\partial_x^{2m}}$, the operator with
cosine-transform multiplier $e^{tu^{2m}}$ (since $\partial_x^{2m}\cos(xu)=(-1)^m
u^{2m}\cos(xu)$): thus $P_{1,t}=e^{-t\partial_x^2}$ is the backward de
Bruijn--Newman flow, $P_{2,t}=e^{+t\partial_x^4}$ the iterated-Brownian flow,
$P_{3,t}=e^{-t\partial_x^6}$, and so on. Then $P_{m,t}$ preserves $\LP$ for all
$t\ge0$ if and only if $m=1$. (The forward Gaussian $e^{+t\partial_x^2}$,
multiplier $e^{-tu^2}$, is the opposite sign and does not preserve $\LP$;
\S\ref{sec:intro}.)
\end{theorem}

\begin{proof}
For $m=1$, $P_{1,t}=e^{-t\partial_x^2}$ preserves $\LP$ by P\'olya's
theorem~\eqref{eq:heatLP}. For each $m\ge2$ an explicit polynomial counterexample fails:
the real-rooted
\[
   f_m(x)=(-1)^m\prod_{j=1}^m\bigl(x^2-j^2\bigr)\in\LP,
\]
of degree $2m$ with leading coefficient $(-1)^m$, has $f_m^{(2m)}=(-1)^m(2m)!$
and $f_m^{(2m+2)}=0$, so $P_{m,t}f_m=f_m+(-1)^m t\,f_m^{(2m)}=f_m+(2m)!\,t$, a
constant upward shift. On $(1,2)$, $f_m<0$ with $f_m(1)=f_m(2)=0$. Since
$f_m<0$ and $x>0$ there, $f_m'=2x\,f_m\sum_{j=1}^m(x^2-j^2)^{-1}$ has the
opposite sign to the sum, which decreases strictly from $+\infty$ to $-\infty$
across $(1,2)$ (each summand has derivative $-2x(x^2-j^2)^{-2}<0$); so $f_m'$ has
a single simple zero $a$, passing from $-$ to $+$. Thus $a$ is a strict minimum,
with $f_m(a)<0$ and $f_m''(a)>0$. For $t$ with $(2m)!\,t>-f_m(a)$, the shift
$g:=P_{m,t}f_m=f_m+(2m)!\,t$ has $g(a)>0$, $g'(a)=0$, and $g''(a)=f_m''(a)>0$: a
strict positive local minimum. Hence $L_1[g](a)=-g(a)\,g''(a)<0$, and
$g\notin\LP$ by Lemma~\ref{lem:laguerre}. For $m=2$ this is the counterexample
$f_2=(x^2-1)(x^2-4)$ of \eqref{eq:witness}.
\end{proof}

\begin{remark}[Universal factors]\label{rem:universal}
By P\'olya's characterization of universal
factors~\cite{Polya27,deBruijn}, a cosine-transform multiplier preserves the
real-rootedness of every admissible transform iff, as a function of $w=u^2$, it
has the form $e^{\gamma w}\prod_k(1+\delta_kw)$ with $\gamma,\delta_k\ge0$. Being
entire and zero-free, $e^{tw^m}$ can carry no $(1+\delta_kw)$ factor, so it must
equal $e^{\gamma w}$, forcing $m=1$. The counterexamples $f_m$ make the failure for
$m\ge2$ concrete, and the order-four case is realized arithmetically by
$H_{8,4}\notin\LP$ (\S\ref{sec:zeta}).
\end{remark}

Thus among the growing-multiplier orders only one, the order-two de Bruijn
flow, preserves real-rootedness, and the first failure is the order-four
iterated-Brownian case. Probabilistically, the Gaussian kernel of Brownian
motion is variation-diminishing (Schoenberg) while the signed biharmonic kernel
of the iterated process is not, so the latter can create real sign changes and,
off the line, non-real zeros; but variation-diminishing is weaker than
$\LP$-preservation: the Gaussian convolution itself does not preserve $\LP$
(\S\ref{sec:intro}), and the order-two preserver is the backward de Bruijn
direction.

%==============================================================
\section{The iterated-Brownian flow of the Riemann \texorpdfstring{$\Xi$}{Xi}}
\label{sec:zeta}
%==============================================================

We turn to the boundary example. With $\Phi$ the Riemann kernel
\[
   \Phi(u)=\sum_{n\ge1}\bigl(2\pi^2n^4e^{9u}-3\pi n^2e^{5u}\bigr)e^{-\pi
   n^2e^{4u}},\qquad
   H_0(x)=\int_0^\infty\Phi(u)\cos(xu)\,du=\tfrac18\Xi(x/2),
\]
the transform $\hat H_0=\Phi$ decays doubly exponentially, so $H_0$ is
admissible and Proposition~\ref{prop:semigroup} applies: the
iterated-Brownian evolution of $H_0$ is the quartic de Bruijn--Newman flow
\[
   H_{t,4}(x)=P_tH_0(x)=\int_0^\infty e^{tu^4}\Phi(u)\cos(xu)\,du .
\]
Define $\Lambda(4)=\inf\{T:H_{t,4}\in\LP\ \text{for all }t\ge T\}$, the quartic
analogue of the de Bruijn--Newman constant.

Since $H_0=\tfrac18\Xi(\cdot/2)$ lies in $\LP$ if and only if the Riemann
hypothesis holds, $H_{t,4}\notin\LP$ at a given $t$ stands on its own, while the
flow's carrying $\Xi$ out of the class, which presumes $\Xi$ began
inside, needs RH.

Under the flow the zeros move. Differentiating $H_{t,4}(x_j(t))=0$ in $t$ and
using $\partial_tH_{t,4}=\partial_x^4H_{t,4}$ gives the velocity law
\[
   \dot x_j(t)=-\,\frac{\partial_x^4H_{t,4}}{\partial_xH_{t,4}}\bigg|_{x=x_j(t)},
\]
so a simple real zero slides along the axis while $\partial_xH_{t,4}\neq0$ and
leaves it only through a collision $\partial_xH_{t,4}=0$ of two real zeros. The
pointwise certificate detects a completed collision; the saddle analysis below is
an asymptotic evaluation of this transport, locating where the motion forces a
collision and where it cannot.

\subsection{Non-membership at \texorpdfstring{$t=8$}{t=8}}
The certificate (Lemma~\ref{lem:laguerre}) applies to the quartic flow directly:
$H_{t,4}\notin\LP$ as soon as one real $x$ has $L_1[H_{t,4}](x)<0$, that is
$H_{t,4}(x)\,H_{t,4}''(x)>H_{t,4}'(x)^2$, three convergent cosine integrals (the
derivatives bring down $-u\sin xu$ and $-u^2\cos xu$ under the same weight
$e^{tu^4}\Phi$) at a point of one's choosing. The floor that moves the relevant
value is, by Proposition~\ref{prop:semigroup},
\begin{equation}\label{eq:floor4}
   H_{t,4}(x)=H_0(x)+t\,H_0''''(x)+O(t^2),\qquad
   H_0''''(x)=\int_0^\infty u^4\,\Phi(u)\cos(xu)\,du,
\end{equation}
the value at $x$ of a cosine transform of the same arithmetic family as $H_0$.

The first close critical-line pair already yields a certificate, of the
negative-maximum kind. The zeros of $H_0=\tfrac18\Xi(\cdot/2)$ at
$x=2\gamma_4=60.85$ and $x=2\gamma_5=65.87$ bound not a dip but a positive
hump: $H_0>0$ on $(60.85,65.87)$, with $H_0(x_*)\approx2.1\times10^{-9}$ at the
midpoint $x_*\approx63.3$. There the floor is negative, $H_0''''(x_*)<0$, so the
flow lowers the hump; the two real zeros approach, collide near $t\approx6.5$,
and the hump sinks beneath the axis. For $t$ just past, $x_*$ is a strict
negative local maximum of $H_{t,4}$ ($H_{t,4}(x_*)<0$, $H_{t,4}'(x_*)=0$,
$H_{t,4}''(x_*)<0$), so $L_1[H_{t,4}](x_*)=-H_{t,4}(x_*)\,H_{t,4}''(x_*)<0$. The Laguerre certificate at a single point makes this rigorous.

\begin{proposition}[Certified: $H_{8,4}\notin\LP$]\label{prop:leave}
At $t=8$ and $x_0=63.2083$, validated interval arithmetic (Arb) gives
\[
   L_1[H_{8,4}](x_0)\in[-2.3966,\,-2.3965]\times10^{-20}\subset(-\infty,0),
\]
with $H_{8,4}(x_0)\in[-1.6415,-1.6414]\times10^{-10}$ and
$H_{8,4}''(x_0)\in[-1.4601,-1.4600]\times10^{-10}$ both negative and
$|H_{8,4}'(x_0)|<3.7\times10^{-13}$. Hence $H_{8,4}\notin\LP$.
\end{proposition}

\begin{proof}
$L_1[H_{8,4}](x_0)=H_{8,4}'(x_0)^2-H_{8,4}(x_0)\,H_{8,4}''(x_0)$ is a finite
combination of the integrals $\int_0^\infty u^k e^{8u^4}\Phi(u)\,\{\cos,\sin\}(x_0u)\,du$,
$k=0,1,2$. Split each at $U=1.6$.

For $u\ge U$ the kernel $\Phi(u)=\sum_{n\ge1}(2\pi^2n^4e^{9u}-3\pi n^2e^{5u})
e^{-\pi n^2e^{4u}}$ satisfies $\Phi(u)\le2.1\,\pi^2e^{9u-\pi e^{4u}}$ (its $n=1$
term, the constant absorbing the rest), so with $|\{\cos,\sin\}|\le1$ and
$u^k\le u^2$ each tail integrand is dominated by
\[
   M(u):=2.1\,\pi^2u^2\,e^{\,8u^4+9u-\pi e^{4u}}.
\]
Here $\dfrac{d}{du}\log M=\dfrac2u+32u^3+9-4\pi e^{4u}\le-1$ on $[U,\infty)$, the
double-exponential term dominating, so $M(u)\le M(U)\,e^{-(u-U)}$ and
\[
   \Bigl|\int_U^\infty\Bigr|\ \le\ \int_U^\infty M(u)\,du\ \le\ M(U)\ <\ 10^{-700}.
\]
Truncating the series at $n=15$ omits a remainder below $10^{-300}$ on $[0,U]$, by
the same $e^{-\pi n^2e^{4u}}$ decay. Both lie far below the enclosure width
$\sim10^{-24}$ and cannot affect the sign.

On $[0,U]$ the integrals are enclosed by Arb's verified
integrator~\cite{Johansson}, whose ball arithmetic with adaptive certified Gauss
quadrature captures the ten orders of oscillatory cancellation in $H\sim10^{-10}$
that a direct interval sum cannot; a bespoke interval Taylor-model integrator,
integrating $\cos(x_0u)$ and $\sin(x_0u)$ against an interval-polynomial envelope
analytically, returns the same enclosure. Working precision, panel counts, and
validation against a closed-form Gaussian integral are in the accompanying
repository.\footnote{\url{https://github.com/QuijoticResearch/quartic-de-bruijn-flow}}
\end{proof}

The certificate is self-contained: it uses no information about the zeros of
$\zeta$, which serve only to locate where to evaluate. The polynomial counterexample of
\S\ref{sec:nonpres} makes the bare non-preservation rigorous for the abstract
semigroup; Proposition~\ref{prop:leave} makes it rigorous for the arithmetic
function $\Xi$ itself, through the drowning of an unusually close low-lying pair.

\subsection{Real-rooting on an expanding window}
The departure is transient on the bulk of the spectrum, and this is a theorem.
While $t$ is small the weight $e^{tu^4}\Phi(u)$ is maximal at the endpoint $u=0$,
so $H_{t,4}$ is a quartic-heat smoothing of $H_0$; near $t\approx150$ the
maximizing saddle of $g(u)=tu^4+\log\Phi(u)$ detaches to
$u_*(t)\approx\tfrac14\log t$, and steepest descent gives
\begin{equation}\label{eq:bulkcos}
   H_{t,4}(x)=e^{C}\sqrt{2\pi}\,w\,e^{-x^2/2|g''|}\cos(xu_*)\,(1+o(1)),
   \qquad C=g(u_*),\ w=|g''|^{-1/2},
\end{equation}
a Gaussian-damped cosine, real-rooted to leading order, with $|g''|=4tu_*^2(4u_*-3)$.

Carried into the complex plane, this controls the zeros. For $z=x+iy$ split
$\cos(zu)=\tfrac12(e^{izu}+e^{-izu})$, so $H_{t,4}(z)=I_+(z)+I_-(z)$ with
$I_\pm(z)=\tfrac12\int_0^\infty e^{g(u)\pm izu}\,du$. Each $I_\pm$ has a saddle at
$u_\pm=u_*\pm iz/|g''|+\cdots$, of value $S_\pm(z)=C\pm izu_*-z^2/2|g''|$, and the
two together give the leading profile
\begin{equation}\label{eq:Acomplex}
   A_t(z)=\sqrt{2\pi}\,w\,e^{C}e^{-z^2/2|g''|}\cos(zu_*),
\end{equation}
the continuation of \eqref{eq:bulkcos} to complex $z$. Off the real axis the two
saddles are of unequal size: $\Re(S_--S_+)=2yu_*$, so for $y\neq0$ one dominates
the other by $e^{2|y|u_*}$, and $A_t$ vanishes only at the real points
$z_m=(m+\tfrac12)\pi/u_*$, all simple.

\begin{lemma}[Saddle scale]\label{lem:scale}
As $t\to\infty$, $u_*\sim\tfrac14\log t$ and $|g^{(k)}(u_*)|\sim 4^k\,tu_*^3$ for
each fixed $k\ge2$. Hence $\rho:=|g'''(u_*)|/|g''|^{3/2}\asymp(tu_*^3)^{-1/2}\to0$
and $|g^{(k)}(u_*)|/|g''|^{k/2}\asymp(tu_*^3)^{1-k/2}$.
\end{lemma}

\begin{lemma}[Uniform two-saddle approximation]\label{lem:twosaddle}
Let $K_t\to\infty$ with $K_t^3\rho\to0$, where $\rho=|g'''(u_*)|/|g''|^{3/2}$. For
all large $t$ and all $z$ with $|z|\le K_t\sqrt{|g''|}$,
\[
   H_{t,4}(z)=A_t(z)+R_t(z),\qquad
   |R_t(z)|\le \eta_t\cdot\tfrac12\sqrt{2\pi}\,w\,e^{C}e^{-\Re z^2/2|g''|}e^{|y|u_*},
\]
the last factor being the modulus of the dominant saddle, with
$\eta_t\le c\,K_t^3\rho+e^{-cC}\to0$ for an absolute constant $c$.
\end{lemma}

\begin{proof}
By symmetry it suffices to treat $I_+(z)=\tfrac12\int_0^\infty e^{\psi(u)}\,du$,
$\psi(u)=g(u)+izu$, the term $I_-$ being identical under $z\mapsto-z$; write
$w=|g''|^{-1/2}$. The saddle $u_+$ solves $\psi'(u_+)=g'(u_+)+iz=0$; from
$g'(u_*)=0$, $g''(u_*)=-|g''|$ and the implicit function theorem,
$u_+=u_*+iz/|g''|+O(|g'''|\,|z|^2/|g''|^3)$, so for $|z|\le K_t\sqrt{|g''|}$,
\[
   |u_+-u_*|\le K_t w\,(1+o(1))=o(1),
\]
since $K_t w=K_t/\sqrt{|g''|}\asymp(tu_*^3)^{-5/12}\to0$; thus $u_+$ stays in the
strip $|\Im u|<\pi/8$ where $g$ is analytic, with value
$\psi(u_+)=S_+(z)=C+izu_*-z^2/2|g''|+O(K_t^3\rho)$. Deform $[0,\infty)$ to the
steepest-descent contour through $u_+$ and split at the neighbourhood
$N=\{\,|u-u_*|\le 2K_t w\,\}$, which contains $u_+$ with margin $K_t w$.

(i) On $N$ the quadratic part of $\psi$ about $u_+$ integrates to
$\tfrac12\sqrt{2\pi/|g''|}\,e^{S_+}=\tfrac12\sqrt{2\pi}\,w\,e^{S_+}$, as
$\psi''(u_+)=g''(u_+)=-|g''|(1+o(1))$.

(ii) The cubic and higher terms give the relative error. The scalings of
Lemma~\ref{lem:scale} hold throughout $N$ (where $e^{4u}$ varies by $1+o(1)$), so
with $\xi=K_t/(tu_*^3)^{1/2}\to0$,
\[
   \sum_{k\ge3}\frac{|g^{(k)}|}{k!}\,(2K_t w)^k
   \;\le\; c\sum_{k\ge3}\frac{tu_*^3}{k!}\,(2\xi)^k\;\le\; c'\,K_t^3\rho,
\]
the factors $4^k$ in $|g^{(k)}|$ and in $|g''|^{k/2}$ cancelling; the classical
steepest-descent remainder~\cite{Olver} then bounds the contribution of $N$ by
$\tfrac12\sqrt{2\pi}\,w\,e^{\Re S_+}\bigl(1+O(K_t^3\rho)\bigr)$.

(iii) The lower endpoint contributes the boundary term
$\tfrac12 e^{g(0)}/|g'(0)+iz|\le\tfrac12 e^{g(0)}/\max(1,|z|)$, which against
$\tfrac12\sqrt{2\pi}\,w\,e^{\Re S_+}$ is $O(e^{g(0)-C})=O(e^{-cC})$ as
$C\to\infty$.

(iv) Off $N$ one has $\Re\psi\le\Re S_+-\tfrac12 K_t^2$ next to the saddle, and
$g(u)\to-\infty$ as $u\to\infty$ by the double-exponential decay
$\Phi(u)\le e^{-\pi e^{4u}}$; this part is $O\!\bigl(e^{\Re S_+}e^{-K_t^2/2}\bigr)$,
negligible against (i).

Summing (i)--(iv), $I_+(z)=\tfrac12\sqrt{2\pi}\,w\,e^{S_+(z)}\bigl(1+O(K_t^3\rho)
+O(e^{-cC})\bigr)$ uniformly for $|z|\le K_t\sqrt{|g''|}$, and likewise $I_-$ with
$S_-$. Their sum is $A_t(z)$ of \eqref{eq:Acomplex}, so
$|R_t(z)|=|I_++I_--A_t|\le\eta_t\cdot\tfrac12\sqrt{2\pi}\,w\,e^{C-\Re z^2/2|g''|}
e^{|y|u_*}$, the dominant saddle, with $\eta_t\le c\,K_t^3\rho+e^{-cC}\to0$ (an
absolute $c$ absorbing the two-saddle sum).
\end{proof}

\begin{theorem}[Real-rooting on an expanding window]\label{thm:window}
For every $\gamma<\tfrac23$ there is $t_0(\gamma)$ such that, for $t\ge t_0(\gamma)$,
every zero of $H_{t,4}$ in the disk $|z|\le|g''|^{\gamma}$ is real and simple. This
window exceeds the zero band $\sqrt{|g''|}$ once $\gamma>\tfrac12$, and widens
without bound.
\end{theorem}

\begin{proof}
Fix $\gamma<\tfrac23$ and set $X_t=|g''|^{\gamma}$. Apply Lemma~\ref{lem:twosaddle}
with $K_t=2|g''|^{\gamma-1/2}=2X_t/\sqrt{|g''|}$: since $\rho\asymp|g''|^{-1/2}$
(Lemma~\ref{lem:scale}), $K_t^3\rho\asymp|g''|^{3\gamma-2}\to0$ as $\gamma<\tfrac23$,
so $\eta_t\to0$. In the strip
$|y|\le h:=\pi/2u_*$ tile by $R_m=[m\pi/u_*,(m+1)\pi/u_*]\times[-h,h]$, each
holding one zero $z_m=(m+\tfrac12)\pi/u_*$ of $A_t$. A rectangle $R_m$ that meets
the disk $|z|\le X_t$ extends at most $2\pi/u_*$ beyond it, hence lies in
$|z|\le X_t+2\pi/u_*<K_t\sqrt{|g''|}$ for large $t$, where Lemma~\ref{lem:twosaddle}
holds. On $\partial R_m$ one has $|\cos(zu_*)|\ge1$:
the vertical sides give $|\cos(zu_*)|=\cosh(yu_*)\ge1$, the horizontal sides
$|\cos(zu_*)|\ge\sinh(hu_*)=\sinh(\tfrac\pi2)>1$. With $|y|\le h$ so
$e^{|y|u_*}\le e^{\pi/2}$,
\[
   \frac{|R_t(z)|}{|A_t(z)|}\le \eta_t\,\frac{\tfrac12 e^{|y|u_*}}{|\cos(zu_*)|}
   \le \tfrac12 e^{\pi/2}\,\eta_t<1
\]
for large $t$, so by Rouch\'e $H_{t,4}$ has a single zero in each $R_m$; as
$H_{t,4}$ is real on $\R$ and $R_m$ is symmetric across it, that zero is real and
simple. For $|y|\ge h$, $|\cos(zu_*)|\ge\sinh(|y|u_*)$ gives
$|R_t|/|A_t|\le\eta_t/(1-e^{-\pi})<1$, so $H_{t,4}=A_t+R_t$ does not vanish there.
Hence every zero of $H_{t,4}$ in $|z|\le X_t$ is one of the real $z_m$.
\end{proof}

So the return is a theorem on a window that widens with $t$: the dip-born pairs
of \S\ref{sec:zeta} cannot survive inside $|z|\le X_t$ once $t$ is large. A scan
to $1.2\sqrt{|g''|}$ at $t=200,300,600,10^3,3\times10^3$, well inside that window,
finds no Laguerre certificate of either kind, corroborating
Theorem~\ref{thm:window}; and the measured $|R_t|/|A_t|$ on the rectangle
boundaries runs from $0.07$ at $t=200$ to $4\times10^{-3}$ at $t=2000$, tracking
the $\rho$ of Lemma~\ref{lem:scale}. The theorem is asymptotic: it holds once
$\eta_t<2e^{-\pi/2}$ in Lemma~\ref{lem:twosaddle}, for $t$ beyond a threshold
$t_0(\gamma)$ we do not make effective, as it carries the constant of the
steepest-descent remainder~\cite{Olver}. At the band scale the margins above
place $t_0$ near order $10^2$; below it the only certified statement is $H_{8,4}\notin\LP$ at the
single point $t=8$ (Proposition~\ref{prop:leave}), the bad set being
numerically present through $t\lesssim150$ (the census) and the band scan clean
from $t=200$, none of this from the theorem.
Tracking that constant would make $t_0(\gamma)$ effective and close the narrative gap.

For $|x|>|g''|^{2/3}$ the leading Gaussian bound no longer applies, because the
saddle moves. The saddle $u_z$ of $I_t(z)$ solves $g'(u_z)=-iz$; as $x=\Re z$
grows it runs from $u_*$ to $\tfrac14\log(x/4\pi)+i\pi/8$, the edge of $\Phi$'s
analyticity strip $|\Im u|<\pi/8$, with $\Re u_z\sim\tfrac14\log x\to\infty$. The
off-axis dominance rate $2\Re u_z$ of Theorem~\ref{thm:window} therefore does not
weaken; while $u_z$ stays inside the strip, extending the two-saddle approximation
is a uniform steepest-descent problem, not an arithmetic one. At the edge
$\Im u_z=\pi/8$ the saddle meets the singularity of $\Phi$ that carries the zeros
of $\zeta$, where the smooth model and the arithmetic coincide and the remainder
is uncontrolled. Far-field exclusion then reduces to the $\Xi$-tail having no
off-line zeros, the lower edge of Proposition~\ref{prop:rh}.

The arithmetic that makes a certificate lives near $u=0$ and is suppressed past
detachment by $e^{g(0)-C}=e^{-C+O(1)}$, $C\to\infty$: at $t=300$ the full transform
and its saddle-window part agree to twenty-two digits. A census finds certificates
for $t\lesssim150$ alone (the first pair near $x\approx63$ at $t\approx8$, later
near $x\approx104,118$ at $t\approx120$, none past $t=150$), and in the tail the
remainder sits below $5\times10^{-2}$ of the leading amplitude at $t=300$, falling
about a decade per decade to $\sim10^{-5}$ at $t=3\times10^5$.

\begin{conjecture}[Finite quartic constant]\label{conj:finite}
$\Lambda(4)<\infty$: $H_{t,4}\in\LP$ for all large $t$, so the set of times at
which $H_{t,4}\notin\LP$ is bounded, comparable to the saddle-detachment scale,
of order $10^2$ in this normalization. (If RH holds, this says the flow
returns $\Xi$ to $\LP$ and keeps it there.)
\end{conjecture}

Theorem~\ref{thm:window} does not reach the far field, and there two readings
stay open which the data choose between without settling: either a non-real pair
migrates outward, staying always beyond $X_t$, so $\Lambda(4)=\infty$, or it is
extinguished and $\Lambda(4)<\infty$. The $e^{-C}$ collapse of the arithmetic and
the widening tail margin favour the second. What would settle it is the same
dominance pushed past the saddle window: that the remainder $B$ stays below the
amplitude of $A$ for $|x|>X_t$ too, for all large $t$. Because $B$ has no definite
sign there, this is a two-sided envelope bound, not the one-sided comparison a
sign-definite floor would permit; it is the quartic counterpart, inequality
reversed, of the monotonicity behind de Bruijn's $\Lambda\le\tfrac12$, and the
gap that remains.

A coarse quadrature suggests that a lifted pair persists
until a Laplace window reaches it at $t\asymp x_*^4$, holding a pair off the axis
at every height and forcing $\Lambda(4)=\infty$; this is an artifact of an
under-resolved saddle. The peak, of width $w\approx0.013$ at the first pair
($(t,u_*)\approx(300,1.38)$), is stepped over by a fixed grid of spacing
$\pi/x\approx0.05$, and the spurious value held that pair off the axis to
$t=300$. Resolved, the pair re-coheres by $t\approx30$, and \eqref{eq:bulkcos}
leaves no replacement past detachment. The pointwise certificate removes the
persistence question altogether: nothing is carried through time.

\begin{remark}[Heat versus iterated Brownian motion]\label{rem:consistency}
The contrast with the order-two flow is, conjecturally, complete. De Bruijn's heat
flow $e^{-t\partial_x^2}$ drives the zeros of $\Xi$ monotonically onto the line
and holds them there ($H_t\in\LP$ for all $t\ge\tfrac12$). The order-four flow does not: $H_{t,4}\notin\LP$ at finite time (certified), and
if RH holds this is an unusually close low-lying pair knocked off the axis, the
saddle-detachment analysis suggesting that as $t$ grows the configuration
relaxes back toward $\LP$. The input is the close pair (in the extreme, the
Lehmer phenomenon~\cite{CSV}), and the detaching saddle is, on present evidence,
what repairs it. Where the heat flow pushes zeros to the line and keeps them, the
order-four flow pushes a close pair off and lets go again, provably on a widening
window and conjecturally throughout.
\end{remark}

\begin{proposition}[Phase diagram of the bad set]\label{prop:rh}
Let $G_4=\{t\ge0:H_{t,4}\in\LP\}$ and $B_4=[0,\infty)\setminus G_4$.
\begin{enumerate}
\item[(i)] $8\in B_4$ (Proposition~\ref{prop:leave}), so $B_4\neq\varnothing$.
\item[(ii)] {\rm(upper edge)} Re-entry is $[T,\infty)\subset G_4$ for some
$T<\infty$, with $T>8$ since $8\in B_4$; Theorem~\ref{thm:window} establishes its
content on the widening window $|z|\le X_t$, Conjecture~\ref{conj:finite}
throughout.
\item[(iii)] {\rm(lower edge)} If $0\in\overline{G_4\cap(0,\infty)}$, that is some
$t_n\downarrow0$ with $t_n\in G_4$ (a fortiori if $[0,\tau]\subset G_4$ for a
$\tau>0$), then $H_0\in\LP$: RH holds.
\end{enumerate}
A proof of RH through this flow thus needs both edges, $B_4\subset[\tau,T]$ with
$0<\tau<T<\infty$. The flow links neither to the other: $e^{t\partial_x^4}$ does
not preserve $\LP$, so $G_4$ is not an up-set and membership at large $t$ does
not transfer downward.
\end{proposition}

\begin{proof}
(i) is Proposition~\ref{prop:leave}, whence $T>8$ in (ii). For (iii), $\LP$ is
closed under locally uniform limits~\cite{Levin} and $H_{t,4}\to H_0$ uniformly
on compacta as $t\downarrow0$ (dominated convergence,
$e^{tu^4}\Phi\le e^{u^4}\Phi\in L^1$ for $t\le1$); so a sequence $t_n\downarrow0$
in $G_4$ forces $H_0\in\LP$, which is RH.
\end{proof}

The edges are not symmetric, and the lower one is strictly stronger than RH:
$H_0\in\LP$ gives only $0\in G_4$, compatible with $\inf B_4=0$ should bad times
accumulate at $0^+$, and does not even yield $[0,\tau]\subset G_4$, since the flow
could leave $\LP$ at once. So the dynamics reach only RH-or-more; pure RH is the
value $0\in G_4$ at $t=0$, not a consequence of the flow. The saddle analysis
controls the upper edge $T$; the lower edge $\tau>0$ lies in the pre-detachment
regime, where the arithmetic is unsuppressed and the two-saddle estimate is
silent. So $\Lambda(4)$ is a re-entry threshold, the upper edge $T$, not the de
Bruijn--Newman constant.

\subsection{The general even order}
The exit-and-return is conjectured to be generic. For each $m\ge2$ consider the
growing-multiplier flow $H_{t,2m}(x)=\int_0^\infty e^{tu^{2m}}\Phi(u)\cos(xu)\,du
=P_{m,t}H_0(x)$, with $P_{m,t}=e^{(-1)^m t\partial_x^{2m}}$; its two steps have
complementary sources. The exit is the non-universality of the multiplier:
by Theorem~\ref{thm:order}, $e^{tu^{2m}}$ is not a universal factor for $m\ge2$,
so $P_{m,t}$ carries some admissible $\LP$ function off the axis; for $\Xi$
the certificate of \S\ref{sec:zeta} gives $H_{8,4}\notin\LP$, while the
computation below indicates $H_{t,6}\notin\LP$ at finite $t$. The
leading floor is now $H_{t,2m}=H_0+(-1)^m t\,H_0^{(2m)}+\cdots$, whose sign,
hence which close pairs drown, depends on $m$, but in each computed case a
drowning pair exists. The return is the arithmetic kernel, and is
independent of $m$: against $\Phi(u)\sim e^{-\pi e^{4u}}$, the weight
$e^{tu^{2m}-\pi e^{4u}}$ has a Laplace saddle that detaches to $u_*(t)$ for any
growing power, and past detachment $C=tu_*^{2m}+\log\Phi(u_*)\to\infty$ suppresses
the arithmetic by $e^{-C}$, leaving the real-rooted cosine \eqref{eq:bulkcos}.
Leave from the multiplier, return from the kernel.

Order six bears this out, and its differences from order four are instructive. The order-six flow $P_{3,t}=e^{-t\partial_x^6}$ (multiplier $e^{tu^6}$)
has $H_{t,6}\notin\LP$ near $t\in[15,30]$, at heights $x\approx87$--$104$,
higher than order four's $x\approx63$ since the floor $-t\,H_0^{(6)}$ carries the
opposite sign and the $u^6$ weight favours higher frequencies, selecting
different pairs. The certificates are of both kinds, a filled dip and a drowned
hump. No certificate is found by $t=300$, where the full transform and its
saddle-window part agree to twenty-six digits ($C\approx1.3\times10^4$): the
arithmetic is suppressed not subtly but astronomically.

\begin{conjecture}[Preservation only at order two]\label{conj:monotone}
Among the growing-multiplier flows $\{P_{m,t}=e^{(-1)^m t\partial_x^{2m}}\}_{m\ge1}$,
order two is the unique one under which the Laguerre--P\'olya class is
forward-invariant: $P_{1,t}=e^{-t\partial_x^2}$ is de Bruijn's flow, which
real-roots $\Xi$ and keeps it real ($H_t\in\LP$, $t\ge\tfrac12$). For every
$m\ge2$, $H_{t,2m}\notin\LP$ at finite $t$ and $H_{t,2m}\in\LP$ for all large
$t$, i.e.\ $\Lambda(2m)<\infty$; Conjecture~\ref{conj:finite} is the case $m=2$.
\end{conjecture}

Because the return is uniform in $m$ (the same saddle detaching
against the same kernel), the dominance estimate of \S\ref{sec:zeta}, proved
once, would give $\Lambda(2m)<\infty$ for all $m\ge2$ together. This would sharpen
Theorem~\ref{thm:order} from preservation to permanence:
conjecturally, order two is the only even order under which $\LP$ is
forward-invariant: the only one that, once it has $\Xi$ in the class, keeps it.

%==============================================================
\section{Concluding remarks}\label{sec:concl}
%==============================================================

The dividing line is the de Bruijn direction.
The backward heat operator $e^{-t\partial_x^2}$, the growing multiplier
$e^{tu^2}$, keeps the Laguerre--P\'olya class invariant; this is the order-two
case and the analytic substrate of de Bruijn--Newman theory. (Forward Gaussian
convolution, the Brownian transition $e^{+t\partial_x^2}$, does not
preserve the class; it is only variation-diminishing, which is weaker.) The
order-four member (the iterated-Brownian, biharmonic semigroup
$e^{t\partial_x^4}$ of Funaki and Allouba--Zheng) does not: the counterexample
$(x^2-1)(x^2-4)$ leaves $\LP$ at $t=3/32$ (\S\ref{sec:nonpres}, where
$e^{t\partial_x^4}f=f+tf''''+\cdots$ fills its dip), and Theorem~\ref{thm:order}
shows no growing-multiplier order past two escapes this.

At the de Bruijn--Newman boundary the same lifting reaches $\Xi$. An unusually
close low-lying pair of $H_0$ ($x=2\gamma_4,2\gamma_5$) brackets a positive hump
over a negative order-four floor; the iterated-Brownian average drowns it near
$t\approx6.5$, and at $t=8$ this is a certified Laguerre violation,
$H_{8,4}\notin\LP$ (Proposition~\ref{prop:leave}); if RH holds, so that $\Xi$
began in $\LP$, the flow has carried it out. The exit is transient on the bulk:
once $t$ passes the saddle-detachment scale, $e^{tu^4}$ concentrates the
transform at $u_*\approx\tfrac14\log t$, and a uniform two-saddle approximation
closed by Rouch\'e proves $H_{t,4}$ real-rooted on a window $|z|\le X_t\to\infty$
(Theorem~\ref{thm:window}). Global return, $\Lambda(4)<\infty$
(Conjecture~\ref{conj:finite}), needs the same dominance carried into the far
field $|x|>X_t$. The two
flows then sit on opposite sides of a divide, though not the one first expected:
de Bruijn's heat flow real-roots the line monotonically and keeps it real,
with threshold in the proven range $0\le\Lambda\le0.22$
(Rodgers--Tao~\cite{RodgersTao}, Polymath15~\cite{Polymath15}); its order-four
counterpart is, conjecturally, non-monotone: outside $\LP$ for a bounded
interval of times, then back.

\subsection*{Further directions}
Several directions remain open. First, characterize the subordinated evolutions
that preserve $\LP$: the order-two and order-one flows do, and whether any other
clock does is open. Second, the far-field estimate behind
Conjectures~\ref{conj:finite} and~\ref{conj:monotone}:
bound the remainder of $\int_0^\infty e^{tu^{2m}}\Phi(u)\cos(xu)\,du$ below the
amplitude of its leading profile for $|x|>X_t$, uniformly in large $t$. The
remainder has no definite sign, so this is a two-sided bound; since the saddle
detaches against the same kernel for every $m$, it would give $\Lambda(2m)<\infty$
for all $m\ge2$, $\Lambda(4)$ included.

\section*{Acknowledgements}
Claude, Gemini, and ChatGPT were used as interactive assistants for drafting,
reorganizing, and generating mathematical arguments and computations, under the
author's direction. The author is responsible for all mathematical content and any
errors. Symbolic verification of the counterexample and the order-dichotomy moment
identities accompanies the source.

%==============================================================
\begin{thebibliography}{99}

\bibitem{AlloubaZheng} H.~Allouba and W.~Zheng, \emph{Brownian-time processes:
the PDE connection and the half-derivative generator}, Ann.\ Probab.\
\textbf{29} (2001), 1780--1795.

\bibitem{BMN} B.~Baeumer, M.~M.~Meerschaert and E.~Nane, \emph{Brownian
subordinators and fractional Cauchy problems}, Trans.\ Amer.\ Math.\ Soc.\
\textbf{361} (2009), 3915--3930.

\bibitem{BorceaBranden} J.~Borcea and P.~Br\"and\'en, \emph{P\'olya--Schur
master theorems for circular domains and their boundaries}, Ann.\ of Math.\
(2) \textbf{170} (2009), 465--492.

\bibitem{Burdzy} K.~Burdzy, \emph{Some path properties of iterated Brownian
motion}, in: Seminar on Stochastic Processes 1992, Progr.\ Probab.\
\textbf{33}, Birkh\"auser, 1993, 67--87.

\bibitem{CSV} G.~Csordas, W.~Smith and R.~S.~Varga, \emph{Lehmer pairs of
zeros, the de Bruijn--Newman constant $\Lambda$, and the Riemann hypothesis},
Constr.\ Approx.\ \textbf{10} (1994), 107--129.

\bibitem{CsordasVarga} G.~Csordas and R.~S.~Varga, \emph{Necessary and
sufficient conditions and the Riemann hypothesis}, Adv.\ in Appl.\ Math.\
\textbf{11} (1990), 328--357.

\bibitem{deBruijn} N.~G.~de Bruijn, \emph{The roots of trigonometric
integrals}, Duke Math.\ J.\ \textbf{17} (1950), 197--226.

\bibitem{Funaki} T.~Funaki, \emph{Probabilistic construction of the solution of
some higher order parabolic differential equation}, Proc.\ Japan Acad.\ Ser.\ A
Math.\ Sci.\ \textbf{55} (1979), 176--179.

\bibitem{Hochberg} K.~J.~Hochberg, \emph{A signed measure on path space related
to Wiener measure}, Ann.\ Probab.\ \textbf{6} (1978), 433--458.

\bibitem{Johansson} F.~Johansson, \emph{Arb: efficient arbitrary-precision
midpoint-radius interval arithmetic}, IEEE Trans.\ Comput.\ \textbf{66} (2017),
1281--1292.

\bibitem{Krylov} V.~Yu.~Krylov, \emph{Some properties of the distribution
corresponding to the equation $\partial u/\partial t=(-1)^{q+1}\partial^{2q}
u/\partial x^{2q}$}, Soviet Math.\ Dokl.\ \textbf{1} (1960), 760--763.

\bibitem{Levin} B.~Ya.~Levin, \emph{Distribution of Zeros of Entire
Functions}, Transl.\ Math.\ Monogr.\ \textbf{5}, Amer.\ Math.\ Soc., 1980.

\bibitem{Nane} E.~Nane, \emph{Higher order PDE's and iterated processes},
Trans.\ Amer.\ Math.\ Soc.\ \textbf{360} (2008), 2681--2692.

\bibitem{Newman} C.~M.~Newman, \emph{Fourier transforms with only real zeros},
Proc.\ Amer.\ Math.\ Soc.\ \textbf{61} (1976), 245--251.

\bibitem{Olver} F.~W.~J.~Olver, \emph{Asymptotics and Special Functions},
Academic Press, New York, 1974.

\bibitem{Polymath15} D.~H.~J.~Polymath, \emph{Effective approximation of heat
flow evolution of the Riemann $\xi$-function, and a new upper bound for the de
Bruijn--Newman constant}, Res.\ Math.\ Sci.\ \textbf{6} (2019), no.~31.

\bibitem{Polya13} G.~P\'olya, \emph{\"Uber Ann\"aherung durch Polynome mit
lauter reellen Wurzeln}, Rend.\ Circ.\ Mat.\ Palermo \textbf{36} (1913),
279--295.

\bibitem{Polya27} G.~P\'olya, \emph{\"Uber trigonometrische Integrale mit nur
reellen Nullstellen}, J.\ Reine Angew.\ Math.\ \textbf{158} (1927), 6--18.

\bibitem{PolyaSchur} G.~P\'olya and J.~Schur, \emph{\"Uber zwei Arten von
Faktorenfolgen in der Theorie der algebraischen Gleichungen}, J.\ Reine
Angew.\ Math.\ \textbf{144} (1914), 89--113.

\bibitem{RodgersTao} B.~Rodgers and T.~Tao, \emph{The de Bruijn--Newman
constant is non-negative}, Forum Math.\ Pi \textbf{8} (2020), e6.

\bibitem{Titchmarsh} E.~C.~Titchmarsh, \emph{The Theory of the Riemann
Zeta-function}, 2nd ed., revised by D.~R.~Heath-Brown, Oxford Univ.\ Press,
1986.

\end{thebibliography}

\end{document}
